\section{방법론}
\label{sec:method}

이 장은 공개 경기 기록에서 예측 표본을 만들고(A, B), 세 연구 질문에 답하는 모델을 세우고(C), 이들을 공정하게 비교하는(D) 과정을 그림~\ref{fig:pipeline}의 순서대로 설명한다.

\begin{figure*}[t]
\centering
\includegraphics[width=\textwidth]{fig_pipeline_clean.pdf}
\caption{데이터 파이프라인과 모델링·평가 절차}
\label{fig:pipeline}
\end{figure*}

%======================================================================
\subsection{데이터 수집}
\label{sec:method:data}
누가 언제 탈락하는지 분석하려면 경기 내내 모든 팀이 어디에 있었고 무엇을 했는지 알아야 한다.
PUBG의 공개 텔레메트리는 경기에 참가한 모든 플레이어의 위치를 약 10초마다 기록하고, 피해, 기절, 처치, 아이템, 안전구역의 상태, 단계 전환 같은 사건을 일어난 시각과 함께 기록한다~\cite{PUBGAPI}.
본 연구는 공개 API가 제공하는 직전 24시간의 표본 경기 목록을 2026년 3월 27일--6월 2일(UTC)에 35회 조회하여, 3월 25일--6월 1일에 생성된 스팀 서버의 경기를 수집하였다.
이 기간의 게임 버전은 PC 업데이트 40.2--41.2이며~\cite{PUBGPatchNotes2026}, 일부 맵 요소를 빼면 안전구역 규칙은 같고 규칙 개편~\cite{PUBGBlueZone2026}은 수집 이후에 적용되었다.
경기 중 모든 플레이어의 위치를 이 정도로 촘촘하게 담은 공개 자료는 드물며, 기존 배틀로얄 예측 연구가 사용한 경기 내 기록도 공개되지 않았다~\cite{Sagong25a,KRAFTON}. 사용한 경기와 표본의 규모는 표~\ref{tab:cohorts}와 같다.

%======================================================================
\subsection{누수 없는 표본 구성}
\label{sec:method:samples}
예측 시점보다 뒤에 일어난 일이 입력에 조금이라도 섞이면, 모델은 결과를 미리 보고 맞히는 셈이 된다.
한 팀의 예로 전체 과정을 먼저 보이면 다음과 같다(그림~\ref{fig:pipeline}의 B 아래 시간축).
어느 경기에서 3번째 안전구역으로의 축소가 경기 시작 후 1,200초에 끝나고 그 순간 4번째 구역이 공개된다고 하자.
그러면 단계 3의 예측 시점은 1,200초, 특징 절단 시각은 1,190초이며, 팀 A의 입력은 1,190초 이전의 기록, 곧 그때까지의 팀원 위치와 생존 여부, 주변 팀과의 거리와 교전, 공개된 구역에 대한 위치만으로 만든다. 1,195초에 일어난 교전은 입력에 들어가지 않는다.
팀 A가 1,200초에 살아 있으면 단계 3의 위험집합에 들어가고, 4번째 구역으로의 축소가 끝나기 전에 마지막 팀원이 처치되면 라벨이 1, 끝까지 살아 있으면 0이다.
이렇게 단계 경계마다 살아 있는 팀 하나가 표본 하나가 되며, 아래에서 각 단계를 정의한다.

\subsubsection{시간축, 예측 시점, 입력 절단}
먼저 모든 기록이 같은 시계를 쓰도록 각 기록의 타임스탬프를 경기 시작부터 흐른 시간(초)으로 바꾼다.
경기 $m$의 단계 $k$는 $k$번째 안전구역이 공개된 때부터 그 구역으로의 축소가 끝날 때까지이다.
축소가 끝나 다음 안전구역이 공개되는 순간, 즉 각 팀이 다음 이동을 정해야 하는 순간을 예측 시점 $b_{m,k}$로 둔다.
입력은 그보다 $\Delta=10$초 앞선 특징 절단 시각 $a_{m,k}=b_{m,k}-\Delta$ 이전의 기록만으로 만든다.
이 규칙은 생존 판정, 팀 간 관계, attention 마스크에도 똑같이 적용된다.
처치된 플레이어는 $a_{m,k}$ 이전에 체력이 남은 위치 기록이 다시 나타날 때만 살아 돌아온 것으로 본다.
무작위로 고른 경기에서 $a_{m,k}$ 이후의 기록을 모두 지우고 입력을 다시 만들어도 원래 입력과 똑같았다.

\subsubsection{사건 기반 라벨과 예측 대상}
정답, 즉 라벨은 “다음 단계가 끝나기 전에 그 팀이 전멸하는가”이다.
기절한 플레이어는 다시 일어날 수 있으므로 전멸로 보지 않으며, 경기 $m$에서 팀 $i$의 탈락 시각 $T_{m,i}$는 다시 돌아오지 않은 마지막 팀원이 처치된 시각이다(우승 팀은 $T_{m,i}=\infty$, 경기 종료 시각 $C_m$까지 관측).
단계 $k$에서 예측의 대상이 되는 팀은 예측 시점에 아직 살아 있는 팀들, 즉 위험집합 $\mathcal{R}_{m,k}=\{i: T_{m,i}>b_{m,k}\}$이다.
각 팀의 라벨은 $Y_{m,i,k}=\mathbb{1}[T_{m,i}\le b_{m,k+1}]$이며, $\mathbb{1}[\cdot]$은 괄호 안이 참이면 1, 거짓이면 0이다. 단계 $k+1$이 경기의 마지막 단계이면 $b_{m,k+1}=C_m$으로 둔다.
표본 $(m,i,k)$는 단계 $k=1,\dots,7$의 경계에서 만들며(솔로 경기에서는 팀이 곧 플레이어), 모델이 추정할 양은 다음 단계 탈락 위험
\begin{equation}
\lambda_{m,i,k}=\Pr\!\left(Y_{m,i,k}=1 \mid i\in\mathcal{R}_{m,k},\,\mathcal{H}_{m}(a_{m,k})\right)
\label{eq:hazard}
\end{equation}
이다. 여기서 $\mathcal{H}_{m}(t)$는 경기 $m$에서 시각 $t$ 이전에 기록된 텔레메트리 전체이다.
이는 기준 시점의 생존자만 놓고 다음 구간의 사건확률을 보는 landmark 방식의 이산시간 hazard이다~\cite{SingerWillett93,vanHouwelingen07}.

\subsubsection{코호트와 분할}
\label{sec:method:cohort}
맵, 시점, 팀 규모가 다르면 이동 거리와 교전 양상이 달라지므로, 같은 조합의 경기끼리 묶어 코호트로 분석한다.
분석 대상은 일반 경기와 랭크 경기이며, \rev{튜토리얼·훈련장 경기, 계정 대부분이 AI인 캐주얼 모드, 비플레이어 전투 계정이 많은 Vikendi 맵은 실제 경쟁과 달라 제외하였다.}
일반 경기에도 AI 계정이 섞여 있으므로 각 팀의 AI 계정 비율을 입력에 넣었다.
경기가 충분한 12개 조합을 코호트로 사용하였고(표~\ref{tab:cohorts}), 같은 경기의 팀들이 학습과 시험에 나뉘어 들어가지 않도록 경기 단위로 학습, 검증, 시험 집합을 8:1:1로 고정 분할하였다.

\begin{table}[t]
\centering
\caption{코호트별 경기 수, 팀--단계 표본 수, 다음 단계 탈락률}
\label{tab:cohorts}
\footnotesize
\setlength{\tabcolsep}{2.7pt}
\begin{tabular}{lccrrrr}
\hline
맵 & 시점 & 인원 & 경기 & 학습/검증/시험 & 표본 & 탈락률(\%) \\
\hline
Erangel & 3인칭 & 1 & 224 & 189/15/20 & 31,740 & 35.1 \\
Erangel & 3인칭 & 2 & 656 & 526/65/65 & 59,170 & 29.5 \\
Taego & 3인칭 & 2 & 573 & 459/57/57 & 47,223 & 29.2 \\
Erangel & 1인칭 & 4 & 430 & 344/43/43 & 21,849 & 27.8 \\
Erangel & 3인칭 & 4 & 3,283 & 2,638/321/324 & 196,442 & 25.7 \\
Karakin & 3인칭 & 4 & 333 & 267/33/33 & 12,669 & 32.1 \\
Miramar & 3인칭 & 4 & 1,543 & 1,235/154/154 & 92,279 & 24.6 \\
Paramo & 3인칭 & 4 & 333 & 267/33/33 & 9,963 & 30.5 \\
Rondo & 3인칭 & 4 & 1,039 & 833/103/103 & 53,683 & 25.2 \\
Sanhok & 3인칭 & 4 & 940 & 752/94/94 & 48,793 & 27.6 \\
Taego & 1인칭 & 4 & 435 & 349/43/43 & 23,390 & 30.0 \\
Taego & 3인칭 & 4 & 4,159 & 3,329/415/415 & 220,565 & 24.9 \\
\hline
합계 & & & 13,948 & 11,188/1,376/1,384 & 817,766 & 26.6 \\
\hline
\end{tabular}
\\[2pt]
\parbox{\columnwidth}{\scriptsize 2026년 3월 25일--6월 1일(UTC) 생성 경기(업데이트 40.2/41.1/41.2: 4,439/6,971/2,538경기)이며, 텔레메트리 기록은 약 3억 2,240만 건(위치 9,643만, 사건 2억 2,597만)이다. 탈락률은 다음 단계 안에 탈락한 표본의 비율이다.}
\end{table}

%======================================================================
\subsection{설계와 분석}
\label{sec:method:design}
팀의 상태를 나타내는 서식지 맥락 표현을 먼저 정의한 뒤, 이를 쓰는 세 분석을 연구 질문 순서로 설명한다(그림~\ref{fig:pipeline}의 C).

\subsubsection{서식지 맥락 표현}
\label{sec:method:representation}
줄어드는 서식지에서 생물 집단의 존속이 위치와 경쟁에 달려 있듯이, 팀의 생존도 구역에 대한 위치와 주변의 붐빔에 달려 있다고 보고 팀의 상태를 네 가지로 나타낸다.
구역 경계에 대한 노출은 팀이 구역 밖으로 얼마나 나가 있는지(구역 밖 깊이) 또는 안쪽으로 얼마나 들어와 있는지(구역 안쪽 깊이)이다.
이동 부담은 공개된 다음 안전구역에 들어가려면 필요한 속도가 팀이 실제로 움직인 속도보다 얼마나 큰지이다~\cite{Nathan08}.
주변 경쟁 밀도는 팀 주변 200\,m 안에 있는 다른 팀의 수이고, 남은 팀원 수는 절단 시각에 살아 있다고 판정된 팀원의 수이다.
\rev{본 연구의 입력은 이 구성개념을 담은 서식지 맥락 블록과 나머지 팀 특징 블록들로 이루어지며, 표~\ref{tab:blocks}는 각 블록의 대표 특징을 두 기존 설계의 입력과 함께 보인다(그림~\ref{fig:pipeline}의 C에는 서식지 맥락 블록만 표시).} 서식지 맥락 블록은 RQ1의 해석 모델과 분해 단계 (a1)에, 나머지 팀 특징은 (a2)에, 관계 그래프는 PTR-Former에 쓰인다.

\begin{table}[t]
\centering
\caption{입력 블록과 대표 특징}
\label{tab:blocks}
\footnotesize
\setlength{\tabcolsep}{3pt}
\renewcommand{\arraystretch}{1.0}
\begin{tabular}{p{1.8cm}p{5.4cm}r}
\hline
블록 & 대표 특징 & 수 \\
\hline
\multicolumn{3}{l}{\emph{본 연구}} \\
\textbf{서식지 맥락} & \textbf{구역 밖·안쪽 깊이, 경계까지 거리, 필요 속도와 이동 부족분, 반경 100--500\,m의 적 팀 수, 구역 크기와 축소율, 구역 진입 혼잡도} & 44 \\
경기 맥락 & 단계 번호, 단계 길이, 팀 규모, 생존 팀 수 & 6 \\
팀 상태 & 생존 인원, 평균·최소 체력, 팀원 간 거리 & 9 \\
이동 & 경로 길이, 평균 속도, 굴곡도, 구역 진입 시점 & 7 \\
교전 & 가한·받은 피해, 처치·기절, 교전 상대 수 & 13 \\
자원 & 아이템 획득·장착·사용, 회복, 차량 이용 & 16 \\
공간·고도 & 가장 가까운 적까지 거리, 적 밀도, 고도 & 12 \\
관계 요약 & 근접·교전·공동 압박 관계의 연결 수와 강도 & 43 \\
관계 그래프 & 근접, 교전(보낸·받은), 공동 압박 가중치 & --- \\
\hline
\multicolumn{3}{l}{\emph{기존 설계}} \\
Sagong 등 & 단계별 공격·처치·피해·어시스트 수, 이동 거리, 밀집도, 아이템, 교전·활동 점수, 플레이 유형, 유형별 생존율 & 12 \\
KRAFTON & 처치·기절, 무기별 피해, 받은 피해와 회복, 장비 수준, 주변 아군·적 수, 블루존·흰 원까지 거리, 팀 체력과 생존 인원, 팀원 위치 & 55+ \\
\hline
\end{tabular}
\end{table}
또한 같은 단계의 팀들 사이에는 가까이 있는 관계, 서로 싸운 관계, 함께 구역의 압박을 받는 관계의 세 가지를 정의한다.

\subsubsection{RQ1: 생태 hazard 모델}
\label{sec:method:rq1}
RQ1은 각 요인이 위험을 어느 방향으로 얼마나 바꾸는지를 묻기 때문에, 계수를 그대로 읽을 수 있는 단순한 모델을 쓴다.
구성개념만으로 위험을 설명하는 로지스틱 회귀
\begin{equation}
\mathrm{logit}\,\lambda_{m,i,k}=\alpha_k+\textstyle\sum_{c\in\mathcal{C}}\beta_c z_{c,m,i,k}
\label{eq:ehm}
\end{equation}
를 코호트마다 추정한다.
여기서 $\mathrm{logit}\,p=\log\{p/(1-p)\}$이고, $\alpha_k$는 단계 $k$의 기본 위험 수준(절편)이다.
$\mathcal{C}$는 구역 밖 깊이, 구역 안쪽 깊이, 이동 부족분, 주변 경쟁 밀도, 남은 팀원 수의 다섯 변수와, 남은 팀 수의 영향을 걸러 내기 위한 통제 변수인 절단 시각의 생존 팀 수의 로그로 이루어지며, 솔로 코호트에서는 남은 팀원 수를 뺀다.
$z_{c,m,i,k}$는 변수 $c$를 그 코호트의 학습 집합 평균과 표준편차로 표준화한 값이고, $\beta_c$는 그 변수가 한 표준편차 늘 때 위험의 로짓이 바뀌는 양이다.
표준오차는 같은 경기의 표본끼리 묶어 계산하고, 12개 코호트의 계수는 확률효과 모형으로 합쳐 방향, 크기, 코호트 간 차이, 부호가 같은 코호트의 수를 보고한다~\cite{DerSimonian86}.
요인의 효과가 단계에 따라 달라지는지는 구성개념과 단계 구간(초반 1--2, 중반 3--5, 후반 6--7)의 상호작용으로 검정한다.

\subsubsection{RQ2: 본 연구 설계와 기존 설계}
\label{sec:method:rq2}
본 연구의 설계는 “다음 단계에 탈락할 확률” 자체를 맞히도록 학습하며, 단계마다 다른 기본 위험을 따로 둔다.
즉 위험을
\begin{equation}
\mathrm{logit}\,\lambda_{m,i,k}=\alpha_k+f(\mathbf{x}_{m,i,k},\mathcal{G}_{m,k})
\label{eq:logit}
\end{equation}
로 둔다. 여기서 $\mathbf{x}_{m,i,k}$는 팀 $i$의 서식지 맥락을 포함한 입력 특징 벡터, $\mathcal{G}_{m,k}$는 $a_{m,k}$에 살아 있는 팀들 사이의 세 관계를 담은 그래프, $f$는 학습되는 함수이다.
모델은 학습 집합의 모든 팀--단계 표본에 대한 교차엔트로피 $-\sum[Y_{m,i,k}\log\lambda_{m,i,k}+(1-Y_{m,i,k})\log(1-\lambda_{m,i,k})]$를 최소화한다.
설계가 특정 학습기에 기대지 않음을 보이기 위해 두 학습기를 쓴다.
하나는 단계 번호를 입력으로 받는 그래디언트 부스팅 결정나무(GBDT)~\cite{Friedman01}이고, 다른 하나는 구조의 기여를 분해하기 위한 참조 모델인 PTR-Former이다.
PTR-Former는 같은 팀의 이전 단계들을 돌아보는 attention(단계 축)과 같은 단계의 다른 생존 팀들을 함께 보는 attention(팀 축)을 번갈아 쓰고~\cite{Vaswani17}, 팀 축 attention에는 관계의 종류에 따라 학습되는 편향을 더한다~\cite{Ying21}\footnote{GBDT는 결정나무 200개와 학습률 0.06을, PTR-Former는 차원 128, attention 헤드 4개, 블록 2개, dropout 0.1, 순위 손실 가중치 0.1, AdamW(학습률 $10^{-3}$)와 코사인 학습률 감소를 사용하였고, 최대 12 epoch 동안 검증 로그 손실을 기준으로 조기 종료하였다.}.

KRAFTON의 예측기는 이와 달리 “누가 먼저 탈락하는가”라는 순서를 배운다.
구체적으로는 Cox 부분우도의 항 $\exp(s_i)/\sum_{j\in\mathcal{R}}\exp(s_j)$에 가중치를 둔 목적을 학습하며~\cite{Cox72}, $s_i$는 한 시점에서 팀 $i$의 점수, $\mathcal{R}$은 그 시점에 살아 있는 팀의 집합이다.
모든 팀의 점수에 같은 값을 더해도 이 항은 그대로이므로, 이 목적은 “단계 4가 단계 2보다 전체적으로 더 위험하다”는 정보, 곧 식~\eqref{eq:logit}의 $\alpha_k$를 배우지 못한다.
그래서 점수에 단계별 기본 위험을 덧붙인 확률 $1-\exp\{-\exp(\gamma_k+s_{m,i,k})\}$도 함께 비교한다. 여기서 $s_{m,i,k}$는 예측 시점 $b_{m,k}$에서 팀 $i$의 점수이고, $\gamma_k$는 위험집합의 관측 탈락과 점수로부터 추정한 단계 $k$의 기본 위험이다~\cite{PrenticeGloeckler78}.

두 기존 설계는 원래의 입력, 학습 목적, 학습 방식을 그대로 두고, 같은 분할과 같은 위험집합에서 평가하였다. 본 연구 설계와의 차이와 적용 방식은 표~\ref{tab:compare}에 정리하였다.
또한 적용판을 우리 과제에 불리하게 만들지 않았는지 확인하기 위해, 각 적용판을 원래 과제의 지표(우승 팀 정확도와 탈락 순서 C-index, 솔로 코호트의 개인 생존 분류)로도 평가하고 본 연구의 설계와 비교한다.

\begin{table*}[t]
\centering
\caption{기존 설계와 본 연구 설계의 비교}
\label{tab:compare}
\footnotesize
\renewcommand{\arraystretch}{1.1}
\begin{tabular}{p{2.3cm}p{4.3cm}p{4.6cm}p{4.6cm}}
\hline
 & Sagong 등~\cite{Sagong25a} & KRAFTON~\cite{KRAFTON} & 본 연구 \\
\hline
분석 단위 & 플레이어 & 팀 & 팀(솔로 경기에서는 개인) \\
예측 대상 & 다음 반 단계의 생존 여부 & 탈락 순서의 상대위험(우승확률로 변환) & \textbf{다음 단계 안에 탈락할 조건부 확률} \\
예측 시점 & 반 단계 경계 & 경기 중 여러 시점 & 단계 경계(다음 안전구역 공개 시점) \\
단계별 위험 수준 & 플레이 유형별 전이 생존율 & 표현하지 않음 & \textbf{단계별 기저위험으로 명시} \\
팀 간 상호작용 & 다루지 않음 & 팀 간 attention & 팀 축·단계 축 attention과 관계 편향 \\
학습 목적 & 생존 분류 & 가중 Cox 부분우도 & \textbf{다음 단계 탈락의 교차엔트로피} \\
\hline
본 연구의 적용 & 2·4인 팀은 팀 단위로 집계, 예측 시점을 단계 경계로 통일 & 입력을 본 연구의 텔레메트리로 재구성, 출력을 확률로 보정 & --- \\
\hline
\end{tabular}
\end{table*}

\subsubsection{RQ3: 단일 요인 분해}
\label{sec:method:decomposition}
RQ3는 성능 차이가 어디서 오는지를 묻는다. 그래서 한 번에 한 요소만 바꾸어 가며 성능이 얼마나 변하는지 비교하고, 각 단계의 변화를 그 요소의 효과로 읽는다.
입력의 효과는 학습기를 GBDT로 고정하고 측정한다. Sagong 등의 입력에 서식지 맥락을 더하고(a1), 나머지 입력을 본 연구의 입력으로 바꾸며(a2), 따로 KRAFTON 입력을 본 연구의 입력으로 바꾼다(b).
학습 목적과 기본 위험의 효과는 KRAFTON의 구조에 본 연구의 입력을 넣고 Cox 목적으로 학습한 모델에서 출발하여 측정한다. 여기에 단계별 기본 위험을 덧붙이고(c1), 학습 목적을 다음 단계 교차엔트로피로 바꾸고(c2), 단계 절편을 명시적으로 더한다(d).
구조의 효과는 KRAFTON 구조를 팀 축만 남긴 PTR-Former로 바꾸고(e), 단계 축과 관계 편향을 더하여(f) 측정하며, 두 단계에서는 순위 보조 손실을 쓰지 않는다.
마지막으로 GBDT에 KRAFTON 입력을 더하여 두 표현이 서로 보완하는지 확인한다(g).
이와 별도로 단계 축과 팀 축을 켜고 끄는 2$\times$2 실험, 관계를 무작위로 섞은 대조 모델, 구역을 별도 토큰으로 넣은 모델로 구성한 다섯 비교는 실험 전에 미리 정해 두었다.

%======================================================================
\subsection{평가}
\label{sec:method:evaluation}
모든 설계는 같은 시험 표본에서 같은 지표로 평가하고, 그림~\ref{fig:pipeline}의 D와 같이 코호트별 차이를 하나로 합친다.

\subsubsection{평가 지표}
주 지표인 AUC는 실제로 탈락한 팀에 탈락하지 않은 팀보다 높은 위험을 줄 확률이다.
단계가 뒤로 갈수록 모든 팀의 위험이 함께 올라 단계만 맞혀도 AUC가 높아질 수 있으므로, 같은 단계끼리만 비교한 단계 안 AUC와 같은 경기·같은 순간의 생존 팀끼리만 비교한 경기--단계 내부 AUC를 함께 본다.
예측확률의 정확도는 로그 손실과 Brier 점수로 평가한다.
단계마다 달라지는 위험 수준을 얼마나 잘 나타내는지는 단계별 보정 오차 $\sum_{k}(n_k/n)\,|\bar{p}_k-\bar{y}_k|$로 평가한다. 여기서 $n_k$는 단계 $k$의 시험 표본 수, $n=\sum_k n_k$, $\bar{p}_k$와 $\bar{y}_k$는 그 표본의 평균 예측확률과 실제 탈락률이다.

\subsubsection{통계적 비교}
두 설계의 차이가 우연인지 보기 위해, 시험 경기를 2,000번 다시 뽑는 경기 단위 쌍대 bootstrap으로 코호트마다 신뢰구간을 구한다.
12개 코호트의 결과는 코호트마다 효과가 다를 수 있음을 허용하는 DerSimonian--Laird 확률효과 모형으로 합치고~\cite{DerSimonian86}, REML 추정과 Hartung--Knapp 구간으로 같은 결론이 나오는지 확인한다~\cite{HartungKnapp01}.
AUC 차이가 0.005보다 작으면 실질적으로 같다고 보고, 이를 기준으로 개선, 저하, 실질적 동등, 불확정으로 판정한다.
미리 정한 다섯 비교에는 여러 번 검정한 효과를 보정하는 Holm 방법을 적용하며~\cite{Holm79}, 나머지 비교는 탐색적 분석으로 보고한다.
끝으로 가장 큰 두 코호트에서 학습 초기화를 바꾸고 학습 예산을 두 배로 늘려 결과가 유지되는지 확인한다.
